\thispagestyle{empty}

\selectlanguage{hebrew}

% Single-spaced, as in the English abstract.
\begingroup
\singlespacing
\setlength{\parskip}{2pt}

% Latin runs inside Hebrew are wrapped in babel-hebrew's \LR, which embeds
% them left-to-right. Without it the engine reverses them: OMOP prints as POMO.

\begin{center}
    \Large
    \thesistitlehe

    \vspace{0.25cm}
    \large
    \thesisauthornamehe

    \vspace{0.25cm}
    \large
    עבודת גמר לתואר מוסמך למדעים

    \vspace{0.1cm}
    \large
    אוניברסיטת בן-גוריון בנגב

    \vspace{0.1cm}
    \large
    \thesisyearhe

    \vspace{0.4cm}
    \Large
    \textbf{תקציר}
\end{center}

היפרגליקמיה בקרב מאושפזים היא תופעה שכיחה ומזיקה, והטיפול בה מוסדר בהנחיות קליניות שמידת היישום שלהן נמדדת לעתים רחוקות בלבד. עבודה זו בוחנת האם ניתן להעריך את מידת ההיענות להנחיות באופן אוטומטי ובקנה מידה לאומי, והאם ניתן לנסח את בעיית חיזוי הסיבוכים לטווח הקצר --- המוצגת בדרך כלל כבעיית סיווג מבוססת מאפיינים על גבי מדידות גולמיות --- כבעיית חיזוי המילה הבאה ברצף זמני, באמצעות מודל מבוסס טרנספורמר הפועל על גבי ידע רפואי נרכש.

ייצוג אחד מחבר בין שתי המטרות. בסיס ידע לטיפול בסוכרת, שנבנה במסגרת עבודה זו, מקודד את הסמנטיקה הקלינית של הטיפול באשפוז ואת נייר העמדה הלאומי הישראלי כתבניות זמניות להערכת איכות. מנוע הפשטה זמנית ממיר רשומה בתקן \LR{OMOP} להפשטות זמניות מבוססות ידע, ובאותו מעבר מחשב עד כמה הטיפול המתועד סטה ממפרט ההנחיה. כל מקטע הפשטה הופך למונח קליני מפורש, ולכן החיזוי וההערכה קוראים רצף אחד.

\emph{המחקר המתודולוגי}, על $57{,}078$ אשפוזים ממאגר \LR{MIMIC-IV}, מבסס את הארכיטקטורה. \LR{\textbf{INTERVenE}}, משפחת מודלים מסוג טרנספורמר הפועלת ישירות על מקטעי ההפשטה, כוללת וריאנט דו-כיווני לחיזוי במעבר יחיד ווריאנט אוטורגרסיבי המחולל מסלולי הפשטה עם עקומות סיכון לכל צעד. המקודד משיג \LR{AUPRC} משוקלל של $0.672$, שיפור של $0.041$ על פני קו הבסיס הנוירוני החזק ביותר עם רווחי סמך שאינם חופפים, ושל $0.085$ על פני צינור סימבולי הצורך את אותו רצף מקטעים בדיוק. שיטת ייחוס מבוססת מקודדים דלילים, המחושבת על המודל הפרוס דרך וו מנותק-גרדיאנט, מתרגמת כל הסבר למונח קליני מפורש. ניסויי אבלציה מפרידים בין מקורות השיפור: כיווץ כל מקטע לאירוע נקודתי פוגע בביצועים פגיעה חדה, ואילו החלפת המקטעים בעלי השמות הקליניים בדיסקרטיזציה אוטומטית מבוססת סטיות תקן אינה פוגעת בהם. מכאן שבקוהורט זה ההפשטה מבוססת הידע היא תשתית לפרשנות, המושגת בעלות כמעט אפסית בביצועים, ולא תנאי הכרחי לחיזוי.

\emph{המחקר הקליני}, על $48{,}354$ אשפוזים בשישה בתי חולים כלליים בישראל, מיישם את אותה מערכת בפרקטיקה השגרתית. רמות ההיענות נבדלות לפי \emph{סוג} הפעולה הנדרשת ולא לפי משקלה הקליני: המלצות המסתפקות בפעולה בדידה אחת בקבלה מיושמות כמעט תמיד, ואילו המלצות הדורשות החלטות אינסולין חוזרות לאורך האשפוז מקבלות ציונים נמוכים עד $0.15$ ולרוב נעדרות כליל. ההיענות עלתה במידה מתונה לאחר הטמעתו של מדד איכות לאומי, עם שונות ניכרת בין בתי החולים שהנתונים אינם מסבירים. לאחר שקלול הסתברות הטיפול, היענות מעל החציון נקשרה בסיכוי נמוך ב-$44\%$ לתמותה תוך $30$ יום, וגם בסיכוי נמוך ב-$82\%$ להיפרגליקמיה, לעומת עלייה של $20\%$ בהיפוגליקמיה קלה וללא עלייה מובהקת בהיפוגליקמיה חמורה --- בדיוק חילופי התועלת הלא-סימטריים שהמלצות אלו נועדו להשיג. ההמלצות שההיענות להן הייתה הנמוכה ביותר נמצאו מבין הקשורות ביותר לתוצאים.

ניסוי מסכם מחבר בין שני החלקים. \LR{INTERVenE} עוברת בהצלחה לקוהורט הקליני ומשיגה $0.527$ לעומת $0.472$ של קו הבסיס החזק ביותר. כאן בסיס הידע אכן תורם ליכולת החיזוי: אותה החלפה בדיסקרטיזציה אוטומטית עולה $0.073$ \LR{AUPRC} ומורידה את הזרוע לרמת קווי הבסיס שעל הנתונים הגולמיים, ולכן ממצא \LR{MIMIC-IV} תוחם את תרומת הידע הקליני במעבר למקור חדש ולא באופן גורף. סולם בן שבע זרועות מריץ כל ייצוג קלט עם פרופיל ההיענות הנמדד כקלט למודל ובלעדיו. הוספתו לא שינתה אף מדד באופן מובהק, אף שניתוח הייחוס מראה כי המודל אכן קורא אותו ומדרג את תבניות ההיענות מבין הגורמים המשפיעים ביותר עבור שבעה מתוך שמונה תוצאים. ההיענות מסכמת את המקטעים שהמודל כבר רואה ואינה מוסיפה עליהם --- אמירה על הייצוג, ולא על השאלה אם יש ערך קליני בהיענות להנחיה, שעליה עונה בנפרד ובחיוב הניתוח משוקלל ההסתברות.

\endgroup
\selectlanguage{english}
